\head{3. Why an accepted history gives the answer}
A row represents a bounded tape segment. Each ordinary cell carries a symbol and the head cell also carries a state. A neighborhood of three cells determines each updated cell. Crossing the right wall removes the head permanently. Conversely, finitely many blank cells suffice for any particular halting run. Two tracks fold a two-sided tape at its origin, so this bounded-row model also simulates Mathlib's two-sided machine.
\para{The finite tests}
Each nonfinal record stores its neighbors, its current color $c$ and its proposed next color $d$. A word automaton checks adjacent records, row boundaries, the finite transition table and exactly one final row, at the end, containing a halted head. It recognizes the initial row as data followed by an end marker, blank padding and a fence. The data appear in reverse block order at the immutable end of the candidate. Injectivity of the fixed-width code, read back from that end, forces them to be precisely the data in the seed. The fast probe requires $v_k=c_k$. The delayed probe requires $v_{k+w}=d_k$ at every nonfinal site and absence of a color at every final site.
\para{Why rows cannot be skipped}
Write the current rows as one color string $C$, prefixing each row with a distinguished separator. Likewise form $N$ from the proposed successors of the nonfinal rows. The delayed comparisons give $\operatorname{drop}_w(C)=N$. If there are $r+1$ rows, $N$ has $r$ separators. For $r>0$ it begins at a separator, so the discarded prefix is exactly the first row with its separator. Splitting the suffix shows that every row is the local update of its immediate predecessor. The case $r=0$ reports the initial configuration itself. Induction now identifies the final row. Conversely, write a genuine halting run in a sufficiently wide rectangle and put its width in the header. Its color track extends the seed suffix. Sections 1 and 2 supply the candidate and a round in which the tests pass.
